\section{Introduction}
\label{sec:rs-intro}

The deployment of autonomous multi-agent systems into consequential environments — agricultural resource management, financial orchestration, clinical decision support, and critical infrastructure — has exposed a structural gap that policy alone cannot close. Conventional architectures can describe what agents do; they cannot architecturally guarantee that every agent's authority traces to a human principal. As agent populations spawn recursively, modify state, and execute decisions with real-world consequences, the question of who authorized each action, through what chain of delegation, and with what scope constraints has become the central unsolved problem in agentic infrastructure.

We introduce the \textbf{Recursive Spawning Protocol} (RSP): a structural mechanism built on the \bxthree{} three-layer accountability architecture that makes the delegation chain from every active agent to a human principal a runtime artifact — immutable, cryptographically enforced, and architecturally verifiable — rather than a forensic reconstruction attempted after the fact. RSP is the second named pillar of the \bxthree{} Framework, complementing Loop Isolation, Spatial Firewall, Sandbox Gate, and the Bailout Protocol. Together, these five pillars constitute the upstream accountability guarantee: that \bxthree{} systems fail upward into human consciousness, never downward into autonomous chaos.

% -------------------------------------------------------
\subsection{The Autonomous Drift Problem}
% -------------------------------------------------------

The failure mode RSP targets is \emph{autonomous drift}: the condition in which an agent continues to operate with no traceable authorization path to a human principal. Drift is not a simple system failure — it is more insidious. The agent continues to function, to spawn children, to modify state, to make decisions, with no external indication that its authority has become disconnected from any accountable principal.

We model a multi-agent system as a rooted directed acyclic graph (DAG) of agent nodes. Each node $n$ carries a \emph{parent pointer} $\alpha(n)$ referencing the node that authorized its provisioning. The root of every authorized chain is a human principal $h \in \mathcal{H}$. A \emph{spawn chain} $C = \langle n_0, n_1, \ldots, n_k \rangle$ satisfies $n_{i-1} = \alpha(n_i)$ for all $1 \leq i \leq k$ and $n_0 \in \mathcal{H}$.

\medskip
\noindent\textbf{Definition 1 (Drift Triad).}
An agent enters the \emph{drift triad} when three conditions hold simultaneously: it is \emph{orphaned} — its parent $\alpha(n)$ has terminated or become unreachable — its delegation chain does not reach any human anchor, and it continues to execute its event loop. Formally:
\begin{equation}
\mathrm{DriftTriad}(n) = \langle \mathrm{orphaned}(n),\ \mathrm{drifted}(n),\ \mathrm{operating}(n) \rangle.
\label{eq:intro-drift-triad}
\end{equation}
An agent is \emph{drifted} if its chain $C(n) = \langle n_0, n_1, \ldots, n \rangle$ does not contain a human principal at its origin — equivalently, $\neg\exists\, h \in \mathcal{H} : n_0 = h$. Drift is heritable: a drifted agent $n$ that spawns a child $c$ propagates the drift condition to all descendants. A single undetected drift event at depth $d$ with branching factor $b$ produces up to $b^d$ drifted agents in the field.

Drift arises from three structural conditions endemic to conventional multi-agent runtimes. \emph{Chain opacity}: $\alpha(n)$ is maintained as a mutable reference — when a parent terminates or is reparented, the child's authorization trace is retroactively modified. \emph{Scope mutation}: $R(n)$ can be expanded at runtime by the agent itself without Purpose Layer authorization — successive small expansions accumulate and the chain's effective operating envelope diverges from its human anchor's intent. \emph{Termination ambiguity}: agent termination is event-driven but not structurally required to propagate to children — a parent may exit cleanly without notifying its children, leaving them orphaned and operational.

These three conditions produce four concrete failure modes. \emph{Silent orphaning}: a parent terminates without notifying its child; the child remains operational and drifted. \emph{Scope creep drift}: a node expands its scope without Purpose Layer authorization; successive expansions migrate the chain's effective envelope far from its human anchor's intent. \emph{Re-parenting drift}: $\alpha(n)$ is redirected by an administrative operation, retroactively obscuring the original authorization trace. \emph{Ghost authority}: a node spawns a child without any authorized delegation chain; the child is drifted by construction. In recursive systems, any one of these failure modes at an intermediate node propagates its unauthorized scope to all descendants — the task tree may superficially resemble a correctly authorized decomposition, but its effective operating envelope is unrecognizable as the product of its human anchor's intent.

% -------------------------------------------------------
\subsection{Why Existing Approaches Are Insufficient}
% -------------------------------------------------------

The conventional architectural response to unbounded agent spawning is the \emph{depth limit}: the system enforces a maximum chain depth $D_{\max}$, rejecting any spawn that would exceed this bound. Depth limits are necessary — they prevent infinite spawn cascades — but they are structurally insufficient to prevent autonomous drift.

Consider a chain $C = \langle n_0, n_1, \ldots, n_k \rangle$ where $k < D_{\max}$. The depth constraint is satisfied by construction. Now apply scope creep drift: each node $n_i$ for $i \geq 1$ successively expands its operating scope $R(n_i)$ beyond $R(n_{i-1})$ without Purpose Layer authorization, recording each expansion as a policy update. After $k$ successive expansions, $R(n_k)$ may be entirely disjoint from $R(n_0)$. The chain depth is $k \leq D_{\max}$. The chain has a human anchor at $n_0$. Yet $n_k$ is drifted by definition: its operating scope is not a descendant refinement of its human anchor's intent.

The drift prevention requirement is formally a conjunction of two constraints:
\begin{equation}
\mathrm{Drift\ Prevention} \iff \bigl(\mathrm{depth}(C) \leq D_{\max}\bigr)\ \land\ \bigl(\forall i:\ R(n_{i+1}) \subseteq R(n_i)\bigr).
\label{eq:intro-drift-prevention}
\end{equation}
Depth limits enforce only the first conjunct. A system that enforces $D_{\max}$ but not the scope refinement constraint is immune to unbounded spawn cascades but remains fully vulnerable to scope creep drift. TTL constraints suffer from the same insufficiency — they bound an agent's lifetime, not its scope within that lifetime. Logging and audit infrastructure can partially reconstruct the chain after the fact, but retroactive reconstruction is categorically different from runtime enforcement. A system that discovers, after an incident, that its agent was drifted has not prevented drift — it has documented it. In consequential environments where unauthorized actions may be irreversible, runtime enforcement is a structural requirement, not an optional enhancement.

The deeper reason depth limits fail is that they address the wrong dimension of the accountability problem. They constrain how many spawn generations can exist; they say nothing about what each generation is authorized to do. Without scope refinement as a structural invariant — enforced not by policy convention but by protocol constraint — chain length is bounded while scope divergence is unbounded.

% -------------------------------------------------------
\subsection{Core Insight: Immutable Parent Pointer Chain}
% -------------------------------------------------------

The Recursive Spawning Protocol's core insight is that autonomous drift is structurally impossible when $\alpha(n)$ is immutable — established once at provisioning and unalterable by any actor under any circumstances.

An immutable parent pointer eliminates the conditions that produce drift by construction. Chain opacity cannot arise: the child's parent pointer always traces to the node that authorized its provisioning, regardless of subsequent events — the delegation chain is a first-class runtime artifact established by the Fact Layer write protocol and hash-chained into the tamper-evident forensic ledger. Scope mutation cannot produce drift without Purpose Layer authorization: the Fact Layer pre-commit hook rejects any spawn event where the proposed child scope is not a descendant refinement of the parent's scope — the subset relationship $R(n_{i+1}) \subseteq R(n_i)$ is not a policy convention but a structural invariant enforced at every spawn event. Orphaning without drift is structurally harmless: when a parent terminates, its children remain traceable to their human anchor through the immutable chain; the termination is recorded; the remaining chain is fully auditable.

The immutable parent pointer is not a storage policy or an access control configuration. It is a write protocol constraint enforced by the Fact Layer at the protocol level — before any write is executed, not after it is observed. Any attempt to overwrite $\alpha(n)$ after provisioning is rejected regardless of the requestor's identity, privileges, or justification. There is no administrative override; there is no privileged actor capable of modifying the chain after the fact. Immutability is structural, not contractual.

RSP is the minimal structural mechanism that achieves this guarantee. It adds three constraints to conventional agentic runtimes: an immutable parent pointer established at provisioning and protected by Fact Layer write protocol; a Purpose Layer authorization check at every spawn event requiring scope refinement justification and spawn necessity declaration; and a Fact Layer pre-commit hook enforcing both the depth bound and the scope subset relationship as structural invariants. Together these constraints make autonomous drift architecturally impossible — not statistically unlikely, not conventionally prevented, but structurally impossible as a matter of protocol enforcement.

% -------------------------------------------------------
\subsection{Paper Roadmap}
% -------------------------------------------------------

The remainder of this paper is organized as follows. Section~\ref{sec:rs-background} situates RSP within the research traditions it draws on — agent lifecycle management, accountability and provenance in distributed systems, fixed versus mutable delegation structures, and hierarchical task decomposition — and establishes the \bxthree{} Framework as the minimal sufficient substrate for RSP's accountability guarantees. Section~\ref{sec:rs-drift} provides the formal characterization of autonomous drift, defines the drift triad and its constituent failure sub-modes, catalogs the four concrete failure modes that arise from the three enabling structural conditions, and proves formally that depth limits alone are insufficient to prevent drift. Section~\ref{sec:rs-protocol} specifies the Recursive Spawning Protocol in full: its Purpose Layer validation gate, its Bounds Engine depth management, its Fact Layer immutable pointer and forensic ledger, and its chain verification and convergence criteria. Section~\ref{sec:rs-integration} maps each RSP component to its governing \bxthree{} layer and demonstrates the structural necessity of the three-layer mapping. Section~\ref{sec:rs-correctness} formally states and proves the drift prevention, chain integrity, and termination correctness properties. Section~\ref{sec:rs-limitations} examines the structural costs of RSP's guarantees and identifies directions for future work. Section~\ref{sec:rs-conclusion} concludes with a summary of RSP's contribution to the accountability of multi-agent AI systems.
